%% Multi-target toy manuscript stored in the author-facing input directory.
%%
%% build.sh defines \ManuscriptTarget before reading this file. The
%% \providecommand keeps neutral as the default for direct LaTeX compilation
%% while allowing that command-line selection to take precedence.
\providecommand{\ManuscriptTarget}{neutral}
\documentclass[times,twoside]{_Templates/MultiTargetManuscript}

\usepackage{url}
% Keep the short fixture's final two-column page visually balanced.
% Remove this package if a real manuscript uses another column-balancing tool.
\usepackage{flushend}

% Optional footer overrides (independent of the selected target):
% \journalname{Custom repository or journal}
% \nojournalname

% Give the surname of the lead author for the running footer.
\leadauthor{Example}

\begin{document}
\raggedbottom
\setlength{\emergencystretch}{2em}

\title{A Reproducible Toy Study of Multi-Target Manuscript Generation}
\shorttitle{Multi-Target Manuscripts}

% Each author receives a separate \author command so authblk can associate
% names and affiliations correctly. \Letter marks a corresponding author.
\author[1]{Avery Example}
\author[1]{Blake Sample}
\author[1,\Letter]{Casey Demonstration}
\affil[1]{Department of Reproducible Examples, Example University, Sample City, USA}

\maketitle

\begin{abstract}
Scientific manuscripts are often copied into separate directory trees when only
repository branding differs. That duplication makes later corrections difficult
to propagate and encourages source drift. We present a toy manuscript that uses
one LaTeX class, one bibliography style, and one content source to create
bioRxiv-, arXiv-, Zenodo-, and neutral-profile PDFs. The profiles change only the
footer label; manuscript content, typography, authorship, citations, and figure
sources remain identical. This small example is intended as a development
fixture and as a starting point for a real manuscript, not as evidence of
compliance with every repository's current submission policy.
\end{abstract}

\begin{keywords}
reproducibility | LaTeX | preprints | document automation | open science
\end{keywords}

\begin{corrauthor}
casey.demonstration\at example.edu
\end{corrauthor}

\section*{Introduction}

Reproducible research depends on preserving both the scientific record and the
process that produced it. Shared conventions for data, metadata, and software
improve the ability of other researchers to inspect and reuse scholarly outputs
\cite{wilkinson2016fair}. Manuscript production deserves the same care. When an
article is copied once for each intended repository, prose and references may be
corrected in one copy but not another. The resulting files look related while
quietly representing different states of the work.

The template tested here treats the repository name as configuration rather than
content. A target profile is selected at build time, while the manuscript is read
from the same \texttt{input/Manuscript.tex} file. The target may also be selected with
a native class option for authors who compile interactively. This separation
keeps repository presentation at the edge of the system and keeps scientific
content in one authoritative source.

\section*{Template design}

The design has three components. First, the manuscript contains all authors,
affiliations, narrative text, captions, and citations. Second, the
\texttt{MultiTargetManuscript} class contains shared typography and four small
target profiles. Third, the build helper passes the requested target to LaTeX
before the class is loaded. The bibliography and source-native figure are shared
by every build (Fig.~\ref{fig:workflow}).

\begin{figure*}[!t]
  \centering
  %% Source-native supporting figure for the toy manuscript.
%% This file intentionally uses only core LaTeX constructs so every target
%% profile consumes the exact same portable figure source.
\setlength{\fboxsep}{7pt}%
\begin{tabular}{c@{\quad$\longrightarrow$\quad}c@{\quad$\longrightarrow$\quad}c}
  \fbox{\parbox[c][1.35cm][c]{0.19\textwidth}{\centering
    One manuscript\\One bibliography}}
  &
  \fbox{\parbox[c][1.35cm][c]{0.19\textwidth}{\centering
    Target selector\\Shared class}}
  &
  \fbox{\parbox[c][1.35cm][c]{0.27\textwidth}{\centering
    arXiv $\vert$ bioRxiv $\vert$ Zenodo\\Neutral (no label)}}
\end{tabular}

  \caption{One source tree generates four presentation profiles. The target
  selector changes only the repository label in the footer. The neutral profile
  suppresses that field and its adjacent separator, so it does not leave an
  empty visual artifact.}
  \label{fig:workflow}
\end{figure*}

The four targets are arXiv, bioRxiv, Zenodo, and neutral. The first three values
produce a visible footer label. The neutral value removes the label entirely.
Existing manual controls are retained:
\texttt{\string\journalname\{...\}} supplies a custom label, and
\texttt{\string\nojournalname} suppresses it. These controls make the template
usable for journal drafts without creating another class file.

\section*{Validation strategy}

Document automation should be tested at more than one level. A successful LaTeX
exit status establishes that the source can be processed, but it does not show
that the correct profile was selected or that the footer is legible. The example
therefore supports four complementary checks:

\begin{enumerate}
  \item compile every supported target from a clean auxiliary directory;
  \item inspect logs for undefined citations, references, and fatal errors;
  \item extract text from each PDF and confirm the expected label behavior; and
  \item render the PDF pages and inspect typography, spacing, and clipping.
\end{enumerate}

This layered approach follows the broader principle that computational outputs
should be accessible, inspectable, and reproducible \cite{mesirov2010accessible}.
It also distinguishes a template fixture from a repository submission test. The
fixture proves that target selection changes the local output as designed. It
does not prove that an external service will accept a future upload, because
external requirements and supported TeX environments can change.

\section*{Illustrative result}

All four outputs are expected to have identical pagination and body text. The
arXiv build displays a stylized arXiv label; the bioRxiv build preserves the
stylized bioRxiv mark from the source template; and the Zenodo build displays a
plain Zenodo label. The neutral build retains the author, date, short title, and
page information but contains no repository name. Conditional separator logic
prevents a doubled bar or unexplained gap when the label is absent.

Because the target selection occurs outside the manuscript content, a correction
to a paragraph, reference, author list, or figure is automatically present in the
next build of every profile. This is the principal maintenance advantage of the
consolidated layout. The source tree is also small enough to package for an
archive or transfer to another LaTeX environment without determining which of
several manuscript copies is authoritative.

\section*{Limitations and use}

Repository labels are presentational metadata. They do not imply endorsement by
arXiv, bioRxiv, or Zenodo, and they do not replace each service's submission form
or metadata requirements. Zenodo is a general-purpose repository rather than a
journal with a mandatory article class, so its profile is deliberately a label
on the shared layout. Authors should reserve or add a DOI separately when their
deposit workflow calls for one.

For a real project, replace the toy prose, authorship, bibliography, and figure
while retaining the directory layout and build interface. Authors can add
ordinary LaTeX packages in the manuscript preamble. Changes to page geometry,
title construction, bibliography formatting, or footer policy belong in the
class and should be documented in its dated revision notes.

\section*{Conclusion}

A target selector is sufficient to replace three nearly identical manuscript
trees with one maintainable source. The resulting repository exposes the choice
clearly, keeps supporting files synchronized, and allows neutral or custom output
without another fork of the class. The example is intentionally small so that
future template changes can be compiled and visually checked quickly.

\begin{acknowledgements}
This manuscript is a fictional development fixture. Names, institutions, and
email addresses are illustrative and do not describe real contributors.
\end{acknowledgements}

\section*{Bibliography}
\begingroup
\setlength{\emergencystretch}{5em}
\sloppy
\bibliography{input/References}
\endgroup

\end{document}
